% \begin{lemma}
% \label{thm: bounded solution}
%     Solutions of MC-O-PI, as defined in \eqref{eq: complete difference inclusion}, converge to the region $\qset_\infty$
%     % defined in \eqref{eq: long term region} 
%     with probability one
%     if the learning rates $(\lrate_k)_{k \geq 0}$ and update distributions $(\update_k)_{k \geq 0}$ satisfy the modified Robbins-Monro conditions (\autoref{asp: stochastic conditions on effective learning rate}).
%     % For any $\epsilon>0$ there exists a time step $K$ when the sequence $(q_\kpz)$ generated by MC-OPI is contained for all $k \geq K$ inside the region
% \end{lemma}
\begin{lemma}
    [Bounded solutions]
    \label{thm: bounded limit sets}
    Let $\polset_\infty$ denote the set of policies that appear infinitely often in a sequence $(\pi_k)_{k \ge 0}$ generated by MC-O-PI (\autoref{def: mcopi}).
    % \begin{align}
    %     \polset_\infty \coloneqq \bigcap_{t = 0}^{\infty} \text{cl} \{\pi_k : k \geq t \} ,
    % \end{align}
    If the learning rates $(\lrate_k)_{k \geq 0}$ and update distributions $(\update_k)_{k \geq 0}$ satisfy the modified Robbins-Monro conditions (\autoref{asp: stochastic conditions on effective learning rate}),
    % Under \autoref{asp: standing} 
    the corresponding sequence $(q_k)_{k \ge 0}$ converges almost surely to the set
    \begin{align*}
        % \qset_\infty 
        % \coloneqq
        \Big\{
        q \in \qset : 
        \min_{\pi \in \polset_\infty} q_\pi(s,a) \leq q(s,a) \leq \max_{\pi \in \polset_\infty} q_\pi(s,a),
        % \forall s \in \sset, \forall a \in \aset
        \forall \, (s,a)
        \Big\} .
    \end{align*}
    Consequently,
    % for every solution $(q_k)$ of MC-O-PI,
    there exists an almost surely finite $c_Q \in [0, \infty)$ such that for all $k \geq 0$
        \begin{align*}
            \| q_k \|_\infty^2 < c_Q .
        \end{align*}
\end{lemma}
\vspace{-0.6cm}
% \begin{proof}
%     Let $(q_k)$ and $(\pol_k)$ be solutions of the difference inclusion \eqref{eq: complete difference inclusion} with initial condition $q_0$, learning rates $(\lrate_k)$, update distributions $(\update_k)$ and noise terms $(w_k)$.
%     At every step, the estimate $q_k$ is updated towards action values $q_{\pol_k}$ that satisfy $q_{\worst{\pol}} \leq q_{\pol_k} \leq q_{\best{\pol}}$.
%     % \begin{align*}
%     %     q_{\worst{\pol}} \leq q_{\pol_k} \leq q_{\best{\pol}} .
%     % \end{align*}
%     Therefore, define the sequences $(x_k)$ and $(y_k)$ as
%     \begin{align}
%         x_\kpo &= x_k + \lrate_k \update_k (q_{\worst{\pol}} - x_k + w_k) ,
%         \\
%         y_\kpo &= y_k + \lrate_k \update_k (q_{\best{\pol}} - y_k + w_k) ,
%     \end{align}
%     and initialised at $x_0 = y_0 = q_0$. At every step we have
%     \begin{align*}
%         x_k \leq q_k \leq y_k .
%     \end{align*}
%     Proposition~4.1 and Example~4.3 in \citep{Bertsekas1996Neuro-DynamicProgramming} show that the sequences $(x_k)$ and $(y_k)$ converge with probability one to $q_{\worst{\pol}}$ and $q_{\best{\pol}}$, respectively, because the noise satisfies conditions \eqref{eq: noise prop exp} and \eqref{eq: noise prop var}.
%     As a result, we conclude that almost surely
%     \begin{align*}
%         q_{\worst{\pol}} = \liminf_{k \to \infty} q_k 
%         \leq q_k 
%         \leq \limsup_{k \to \infty} q_k = q_{\best{\pol}} ,
%     \end{align*}
%     % As a result, for every $\epsilon>0$ there exists a time step $K$ such that 
%     % \begin{align*}
%     %     q_{\worst{\pol}} - \epsilon \onefun \leq x_k \leq q_k \leq y_k \leq q_{\best{\pol}} + \epsilon \onefun
%     % \end{align*}
%     % for all $k \geq K$ with probability one.
%     which implies that solutions converge to $\qset_\infty$ with probability one.
% \end{proof}
\begin{proof}
    Let $(q_k)_{k \ge 0}$ and $(\pol_k)_{k \ge 0}$ be solutions of the difference inclusion in \eqref{eq: complete difference inclusion} with initial condition $q_0$, learning rates $(\lrate_k)_{k \ge 0}$, update distributions $(\update_k)_{k \ge 0}$ and stochastic perturbations $(w_k)_{k \ge 0}$.
    
    By assumption, there exists a finite time $K$ such that
    the sequence $(\pol_k)_{k \ge K}$ only contains policies in $\polset_\infty$.
    % the estimate $q_k$ is updated towards action values $q_{\pi_k}$ where $\pi_k \in \polset_\infty$.
    Thus for all $k \ge K$, each estimate $q_k(s,a)$ is updated towards a value $q_{\pol_k}(s,a)$ that satisfies
    \begin{align}
        \min_{\pi \in \polset_\infty} q_\pi(s,a) \leq q_{\pi_k}(s,a) \leq \max_{\pi \in \polset_\infty} q_\pi(s,a) .
    \end{align}
    
    Therefore, define the scalar sequences $(x_k)_{k \ge K}$ and $(y_k)_{k \ge K}$ as
    \begin{align}
        x_\kpo &= x_k + \lrate_k \update_k(s,a) 
        \big( \min_{\pi \in \polset_\infty} q_\pi(s,a) - x_k + w_k(s,a) \big) ,
        \\
        y_\kpo &= y_k + \lrate_k \update_k(s,a) \big( \max_{\pi \in \polset_\infty} q_\pi(s,a) - y_k + w_k(s,a) \big) ,
    \end{align}
    and initialised at $x_K = y_K = q_K(s,a)$. From time step $K$ onwards we have
    \begin{align*}
        x_k \leq q_k(s,a) \leq y_k .
    \end{align*}
    Proposition~4.1 and Example~4.3 in \citep{Bertsekas1996Neuro-DynamicProgramming} show that $x_k$ and $y_k$ converge almost surely to $\min_{\pi \in \polset_\infty} q_\pi(s,a)$ and $\max_{\pi \in \polset_\infty} q_\pi(s,a)$, respectively, because the noise satisfies conditions \eqref{eq: muw_mds} and \eqref{eq: muw_second_moment}.
    As a result, we conclude that almost surely
    \begin{align*}
        &\liminf_{k \to \infty} q_k(s,a)
        \ge \min_{\pi \in \polset_\infty} q_\pi(s,a)
        % \leq q_k(s,a) 
        \\
        &\limsup_{k \to \infty} q_k(s,a)
        \le \max_{\pi \in \polset_\infty} q_\pi(s,a)
    \end{align*}
    for every state-action pair,
    which proves the first part of the lemma.

    Further,
    given that there are a finite number of deterministic policies, $q_k$ converges almost surely to a bounded set.
    As a result, 
    % that sequence is almost surely bounded, meaning that 
    its norm over all $k \geq 0$ is almost surely finite. So there exists an almost surely finite constant $c_Q$ that is strictly greater than $\sup_{k \geq 0} \|q_k\|_\infty^2$.
\end{proof}